\documentclass[blue]{beamer}
\mode<presentation> {
  \setbeamertemplate{background canvas}[vertical shading][bottom=red!10,top=blue!10]
  \usetheme{Boadilla}
  \usefonttheme[onlysmall]{serif}
}
\usepackage{pgf,pgfarrows,pgfnodes,pgfautomata,pgfheaps,pgfshade}
\usepackage{amsmath,amssymb,bm}
\usepackage{colortbl}
\usepackage[english]{babel}
\setbeamercovered{dynamic}
\DeclareMathOperator\PV{PV}
\def\R{\mathbb{R}}
\def\Q{\mathbb{Q}}
\def\N{\mathbb{N}}
\newcommand{\ps}[2]{ \langle #1 , #2 \rangle }
\def\1{\mathbf{1}}
\newcommand{\nor}[1]{\left\| #1 \right\|}
\def\leq{\leqslant}
\def\geq{\geqslant}
\def\et{\quad \textrm{and} \quad}
\def\ou{\quad \textrm{or} \quad}
\def\boite{\hskip 0.5mm {\Box} \hskip 0.5mm}
\DeclareMathOperator\PPE{PPE}
\DeclareMathOperator\vNM{vNM}
\DeclareMathOperator\argmax{argmax}
\DeclareMathOperator\Prob{Prob}
\DeclareMathOperator\supp{supp}
\DeclareMathOperator\Eq{Eq}
\DeclareMathOperator\inte{int}
\DeclareMathOperator\co{co}
\DeclareMathOperator\cl{cl}
\DeclareMathOperator\Dirac{Dirac}
\DeclareMathOperator\F{F}
\DeclareMathOperator\As{As}
\DeclareMathOperator\Ir{Ir}
\DeclareMathOperator\MRS{MRS}
\DeclareMathOperator\Range{Im}
\DeclareMathOperator\DIV{DIV}
%------------------------------ theorem styles
\theoremstyle{definition}
\newtheorem*{assumption}{Assumption}
\newtheorem*{conjecture}{Conjecture}
\theoremstyle{example}
\newtheorem*{proposition}{Proposition}
\newtheorem*{theoremns}{Theorem}
\newtheorem*{remark}{Remark}
\newtheorem*{definitionns}{Definition}
%-----------------------------------------
\title[Corporate Finance]{Theory of Corporate Finance}
\author[Victor Filipe]{V. Filipe Martins-da-Rocha}
\institute[FGV EESP]{Escola de Economia de S\~ao Paulo}
\date[April--June, 2026]{Part 1.c. Stock Market Equilibrium: About competitive price perceptions}
\subject{Economic Theory}
\begin{document}
\frame{\titlepage}

%--------------------------------------------------------------------
\begin{frame}\frametitle{Outline}
\begin{itemize}
\item The Bisin--Gottardi--Ruta framework
\item Equilibrium pricing under short-sale constraints
\item Makowski's competitive price perceptions
\item Strong unanimity
\item Constrained Pareto efficiency
\item Modigliani--Miller revisited
\item Intermediated short sales
\end{itemize}
\end{frame}

%====================================================================
\section{The Framework}
%====================================================================

\begin{frame}\frametitle{Why short-sale constraints?}
Many models assume symmetric trade: any security issued by the firm can be freely bought or short-sold by any investor
\bigskip

This symmetry underlies the Stiglitz homemade-leverage argument
\bigskip

\textbf{In real-world markets, however:}
\begin{itemize}
\item short selling equity requires borrowing shares and posting collateral
\item short positions in corporate bonds are rare
\end{itemize}
\bigskip

When investors are prevented from taking short positions, marginal valuations of a given cash flow may differ across investors at equilibrium, and the Modigliani--Miller logic must be re-examined
\end{frame}

%--------------------------------------------------------------------
\begin{frame}\frametitle{Markets and assets}
\begin{itemize}
\item One firm $J = \{j\}$ (or a continuum of identical firms)
\item One riskless financial asset $A = \{a_0\}$ with $\kappa_s(a_0) = 1$ for all $s \in S$
\item Risk-free price $q$ at $t = 0$, with $q = 1/(1 + r_0)$
\end{itemize}
\bigskip

\textbf{Short-sale constraints:}
\begin{itemize}
\item the firm can issue debt but cannot buy the bond: $\beta \geq 0$
\item investors can buy the bond but cannot issue debt: $\theta^i \geq 0$
\item investors cannot short-sell the firm's equity: $\alpha^i \geq 0$
\end{itemize}
\end{frame}

%--------------------------------------------------------------------
\begin{frame}\frametitle{The firm's technology}
\begin{block}{Production set}
\begin{equation*}
Y := \bigl\{ (y_0, y_1) \in \R \times \R^S \colon y_0 \leq 0 \et y_s \leq f_s(-y_0) \text{ for all $s \in S$} \bigr\}
\end{equation*}
\end{block}
\bigskip

with $f_s : [0, \infty) \to [0, \infty)$ for each $s \in S$
\bigskip

\textbf{Notation:}
\begin{itemize}
\item $k := -y_0 \geq 0$ is the capital invested at $t = 0$
\item $f_s(k)$ is the state-contingent output at $t = 1$
\end{itemize}
\end{frame}

%--------------------------------------------------------------------
\begin{frame}\frametitle{The no-default constraint}
\begin{assumption}[No default]
The firm chooses a capital $k \geq 0$ and a debt level $\beta \geq 0$ such that
\begin{equation*}
\inf \{ f_s(k) \colon s \in S\} \geq \beta.
\end{equation*}
\end{assumption}
\bigskip

The set of \emph{admissible} capital--budgeting and financial strategies is
\begin{equation*}
\mathcal{A} := \bigl\{ (k, \beta) \in \R_+^2 \colon f_s(k) \geq \beta, \quad \text{for all $s\in S$} \bigr\}
\end{equation*}
\end{frame}

%--------------------------------------------------------------------
\begin{frame}\frametitle{Investors and budget sets}
Each investor $i$ takes as given:
\begin{itemize}
\item the firm's capital--budgeting and financial strategy $(k, \beta) \in \mathcal{A}$
\item the security price $q$ and the equity price $E$
\end{itemize}
\bigskip

He chooses
\begin{itemize}
\item a consumption plan $c = (c_0, c_1) \in \R_+ \times \R_+^S$
\item a bond holding $\theta^i \geq 0$
\item a share $\alpha^i \geq 0$
\end{itemize}
\end{frame}

%--------------------------------------------------------------------
\begin{frame}\frametitle{Budget constraints}
\begin{block}{At $t = 0$}
\begin{equation*}
c^i_0 + q \theta^i + E \alpha^i \leq \omega^i_0 + [E + q \beta - k] \xi^i
\end{equation*}
\end{block}

\begin{block}{At $t = 1$ in state $s$}
\begin{equation*}
c^i_s \leq \omega^i_s + \theta^i + [f_s(k) - \beta] \alpha^i
\end{equation*}
\end{block}
\bigskip

The set of admissible actions $(c, \theta, \alpha)$ is denoted $B^i(q, E, k, \beta)$
\end{frame}

%--------------------------------------------------------------------
\begin{frame}\frametitle{Competitive equilibrium given $(k, \beta)$}
\begin{definitionns}
Fix $(k, \beta) \in \mathcal{A}$. A list
\begin{equation*}
\bigl\{ (k, \beta), (q, E), (c^i, \theta^i, \alpha^i)_{i \in I} \bigr\}
\end{equation*}
is a competitive equilibrium if:
\end{definitionns}

\begin{itemize}
\item[(a)] each investor optimizes:
\begin{equation*}
(c^i, \theta^i, \alpha^i) \in \argmax \{ U^i(c) \colon (c, \theta, \alpha) \in B^i(q, E, k, \beta) \}
\end{equation*}
\item[(b)] financial markets clear:
\begin{equation*}
\sum_{i \in I} \theta^i = \beta \et \sum_{i \in I} \alpha^i = 1
\end{equation*}
\item[(c)] consumption markets clear
\end{itemize}
\end{frame}

%====================================================================
\section{Equilibrium Pricing}
%====================================================================

\begin{frame}\frametitle{Present-value operator}
For each investor $i$ with interior consumption $c^i \gg 0$, define
\begin{equation*}
\PV^i[w] := \sum_{s \in S} \MRS^i_s(c^i) w_s, \quad w \in \R^S
\end{equation*}
where $\MRS^i_s(c^i) = \beta^i \nu^i_s (u^i)'(c^i_s) / (u^i)'(c^i_0)$
\bigskip

$\PV^i[w]$ is the value at $t = 0$ of the state-contingent claim $w$, evaluated at investor $i$'s marginal rates of substitution
\bigskip

\textbf{Key fact:} under short-sale constraints, different investors may assign different present values to the same claim
\end{frame}

%--------------------------------------------------------------------
\begin{frame}\frametitle{Max-valuation pricing: theorem}
\begin{proposition}[Max-valuation pricing]
Under \textup{(C.1)--(C.4)}, at any competitive equilibrium with $\beta > 0$:
\begin{equation*}
q = \max_{i \in I} \PV^i[\1_S]
\end{equation*}
\begin{equation*}
E = \max_{i \in I} \PV^i[f_1(k) - \beta \1_S]
\end{equation*}
\end{proposition}
\bigskip

The price of a cash flow is determined by the \emph{maximum} marginal valuation across investors
\end{frame}

%--------------------------------------------------------------------
\begin{frame}\frametitle{Max-valuation pricing: intuition}
The first-order conditions of investor $i$'s problem are
\begin{equation*}
q = \PV^i[\1_S] + \frac{\eta^i_\theta}{\mu^i_0}
\end{equation*}
\begin{equation*}
E = \PV^i[f_1(k) - \beta \1_S] + \frac{\eta^i_\alpha}{\mu^i_0}
\end{equation*}
where $\eta^i_\theta, \eta^i_\alpha \geq 0$ are multipliers on the short-sale constraints, and $\mu^i_0 > 0$ is the Lagrange multiplier on the date-$0$ budget constraint
\bigskip

\begin{itemize}
\item For every investor: $q \geq \PV^i[\1_S]$ and $E \geq \PV^i[f_1(k) - \beta \1_S]$
\item Equality holds when the corresponding short-sale constraint does not bind
\end{itemize}
\bigskip

Market clearing implies at least one binding investor in each market
\end{frame}

%--------------------------------------------------------------------
\begin{frame}\frametitle{A fundamental departure from Arrow--Debreu}
Under complete markets without short-sale constraints:
\begin{itemize}
\item the common state-price deflator equals every investor's MRS
\item all investors are marginal in the market
\end{itemize}
\bigskip

Under short-sale constraints:
\begin{itemize}
\item only the \emph{highest-valuation} investors are marginal
\item all others are strictly infra-marginal
\end{itemize}
\bigskip

The price of a cash flow is determined by the maximum valuation present in the economy
\end{frame}

%====================================================================
\section{Makowski's Price Perceptions}
%====================================================================

\begin{frame}\frametitle{Pricing out-of-equilibrium choices}
Max-valuation pricing pins down equilibrium prices given the firm's actual decision $(k, \beta)$
\bigskip

\textbf{To decide whether $(k, \beta)$ is optimal,} the manager (and each shareholder) must conjecture how prices would respond to an alternative choice $(k', \beta') \in \mathcal{A}$
\bigskip

The bond price $q$ does not depend on $(k, \beta)$
\bigskip

The equity price $E$ depends on $(k, \beta)$ through the firm's dividend stream $(f_s(k) - \beta)_{s \in S}$
\end{frame}

%--------------------------------------------------------------------
\begin{frame}\frametitle{Why Grossman--Hart no longer suffices}
In the previous parts we assumed Grossman--Hart \emph{competitive price perceptions}: each investor evaluates alternative cash flows at his own MRS
\bigskip

Under complete markets, all MRSs coincide, so this gives a single equity-price map
\bigskip

\textbf{Under short-sale constraints:}
\begin{itemize}
\item different investors assign different $\PV^i$'s to the same stream
\item Grossman--Hart no longer yields a single, operational equity-price map
\end{itemize}
\bigskip

We need an alternative criterion
\end{frame}

%--------------------------------------------------------------------
\begin{frame}\frametitle{The Makowski criterion}
\begin{definitionns}[Makowski price perception, Makowski 1983]
Each investor conjectures that the equity price associated with $(k', \beta')$ is the highest marginal valuation that the economy would assign to the resulting dividend stream:
\begin{equation*}
\widetilde{E}(k', \beta') := \max_{i \in I} \PV^i[f_1(k') - \beta' \1_S]
\end{equation*}
for $(k', \beta') \in \mathcal{A}$.
\end{definitionns}
\bigskip

\textbf{Consistency:} at the actual equilibrium choice $(k, \beta)$,
\begin{equation*}
\widetilde{E}(k, \beta) = E
\end{equation*}
by max-valuation pricing
\end{frame}

%--------------------------------------------------------------------
\begin{frame}\frametitle{Makowski: market-microstructure justification}
When a shareholder evaluates the firm's cash flow, he thinks of the benefits he could obtain by selling the firm in the market
\bigskip

He uses the MRS of those investors willing to pay the most --- the \emph{highest-valuation} investors --- because they determine the resale price
\bigskip

In a competitive economy without bilateral matching, selling at a price below $\max_i \PV^i[w]$ would immediately be arbitraged: the highest-valuation buyer would purchase it instead of his current marginal trade
\end{frame}

%--------------------------------------------------------------------
\begin{frame}\frametitle{Makowski: continuum-of-firms interpretation}
Alternative interpretation: a continuum of identical firms, each of measure zero
\bigskip

\begin{itemize}
\item A deviation by a single firm has negligible impact on aggregate quantities
\item Hence on investors' marginal rates of substitution
\item Prices $(q, \widetilde{E})$ are taken as \emph{price maps}, given to firms and investors alike
\end{itemize}
\bigskip

This is analogous to the price-taking assumption of Arrow--Debreu, adapted to the situation in which the relevant ``price'' of a firm is a function of the firm's characteristics
\end{frame}

%--------------------------------------------------------------------
\begin{frame}\frametitle{Makowski: an informational caveat}
\begin{block}{Problem}
Each agent must know (or be able to compute) the MRS of every other investor, or at least the maximum over the population
\end{block}
\bigskip

This is a strong epistemic hypothesis
\bigskip

The literature has addressed it in various ways:
\begin{itemize}
\item perfect competition with a continuum of agents
\item rational-expectations frameworks
\end{itemize}
\bigskip

See Bisin, Gottardi and Ruta (2014) for a discussion
\end{frame}

%====================================================================
\section{Strong Unanimity}
%====================================================================

\begin{frame}\frametitle{The firm's value map}
Under Makowski's criterion, the firm's value is
\begin{equation*}
V(k', \beta') := \widetilde{E}(k', \beta') + q \beta' - k', \quad (k', \beta') \in \mathcal{A}
\end{equation*}
\bigskip

\textbf{Question:} If the firm chooses $(k, \beta)$ to maximize $V$ over $\mathcal{A}$, does every shareholder agree?
\bigskip

The next theorem establishes \emph{strong unanimity}: every shareholder agrees, even when allowed to jointly choose his own consumption and portfolio
\end{frame}

%--------------------------------------------------------------------
\begin{frame}\frametitle{Strong unanimity: theorem}
\begin{theoremns}[Strong unanimity]
Under \textup{(C.1)--(C.4)}, fix a competitive equilibrium
\begin{equation*}
\bigl\{ (k, \beta), (q, E), (c^i, \theta^i, \alpha^i)_{i \in I} \bigr\}
\end{equation*}
in which the firm chose $(k, \beta) \in \mathcal{A}$ to maximize $V$ over $\mathcal{A}$. Then for every $i \in I$:
\begin{equation*}
U^i(c^i) = \max \bigl\{ U^i(c') \colon (c', \theta', \alpha') \in B^i(q, E', k', \beta'), \ (k', \beta') \in \mathcal{A} \bigr\}
\end{equation*}
where $E' = \widetilde{E}(k', \beta')$.
\end{theoremns}
\end{frame}

%--------------------------------------------------------------------
\begin{frame}\frametitle{Strong unanimity: comments}
\textbf{Unanimity:} Every shareholder agrees that the manager should maximize $(k, \beta) \mapsto \widetilde{E}(k, \beta) + q \beta - k$
\bigskip

\textbf{Why ``strong''?} When investor $i$ considers possible deviations $(k', \beta')$ of the manager, he is also allowed to optimize his own action $(c', \theta', \alpha')$ in response
\bigskip

The equilibrium plan is utility-maximizing for every investor over the joint choice of:
\begin{itemize}
\item consumption $c'$
\item bond holding $\theta'$
\item share $\alpha'$
\item firm decision $(k', \beta')$
\end{itemize}
\end{frame}

%====================================================================
\section{Constrained Pareto Efficiency}
%====================================================================

\begin{frame}\frametitle{Admissible allocations}
\begin{definitionns}
A profile $(c^i)_{i \in I}$ is \emph{admissible} if:
\end{definitionns}
\bigskip

\begin{itemize}
\item[(i)] feasibility: there exists $k \geq 0$ such that
\begin{equation*}
k + \sum_{i \in I} c^i_0 \leq \omega_0 \et \sum_{i \in I} c^i_s \leq f_s(k) + \sum_{i \in I} \omega^i_s
\end{equation*}
\item[(ii)] implementability: there exists $\beta \geq 0$ and, for every $i \in I$, $(\theta^i, \alpha^i) \in \R_+^2$ with
\begin{equation*}
c^i_s = \omega^i_s + \theta^i + [f_s(k) - \beta] \alpha^i, \quad \text{for all $s \in S$}
\end{equation*}
\end{itemize}
\end{frame}

%--------------------------------------------------------------------
\begin{frame}\frametitle{Admissibility: comments}
\textbf{Two distinct constraints:}
\begin{itemize}
\item Feasibility: aggregate consumption cannot exceed aggregate endowment plus production
\item Implementability: investor $i$'s consumption must be attainable from non-negative bond and share holdings
\end{itemize}
\bigskip

\textbf{What admissibility does \emph{not} require:}
\begin{itemize}
\item the financial portfolios $(\theta^i, \alpha^i)$ to clear in aggregate
\item a ``planner'' may assign portfolios that do not sum to $\beta$ or to $1$
\end{itemize}
\bigskip

This added flexibility is what makes the constrained-efficiency notion weaker than full Pareto efficiency
\end{frame}

%--------------------------------------------------------------------
\begin{frame}\frametitle{Constrained Pareto efficiency}
\begin{definitionns}
An admissible allocation $(c^i)_{i \in I}$ is \emph{constrained Pareto efficient} if there does not exist another admissible allocation $(\widetilde{c}^i)_{i \in I}$ with
\begin{equation*}
U^i(\widetilde{c}^i) \geq U^i(c^i), \quad \text{for all $i \in I$}
\end{equation*}
and strict inequality for at least one $i$.
\end{definitionns}

\begin{theoremns}[Welfare theorem]
Under \textup{(C.1)--(C.4)}, every competitive-equilibrium allocation is constrained Pareto efficient.
\end{theoremns}
\end{frame}

%====================================================================
\section{Modigliani--Miller Revisited}
%====================================================================

\begin{frame}\frametitle{Equity holders and bond holders}
Fix a competitive equilibrium with $\beta > 0$
\bigskip

\begin{definitionns}[Marginal holders]
\begin{equation*}
I^e := \bigl\{ i \in I \colon E = \PV^i[f_1(k) - \beta \1_S] \bigr\}
\end{equation*}
\begin{equation*}
I^b := \bigl\{ i \in I \colon q = \PV^i[\1_S] \bigr\}
\end{equation*}
\end{definitionns}
\bigskip

\begin{itemize}
\item $i \in I^e$: \emph{equity holder} (on the verge of holding equity)
\item $i \in I^b$: \emph{bond holder} (on the verge of holding the bond)
\end{itemize}
\bigskip

By max-valuation pricing, $I^e$ and $I^b$ are both non-empty
\end{frame}

%--------------------------------------------------------------------
\begin{frame}\frametitle{Equity holders are bond holders}
\begin{proposition}[Equity holders are bond holders]
At a value-maximizing equilibrium with strictly slack no-default constraint
\begin{equation*}
f_s(k) > \beta > 0, \quad \text{for all $s \in S$}
\end{equation*}
we have $I^e \subset I^b$, i.e.,
\begin{equation*}
E = \PV^i[f_1(k) - \beta \1_S] \;\Longrightarrow\; q = \PV^i[\1_S]
\end{equation*}
\end{proposition}
\bigskip

\textbf{Sketch:} differentiating $V(k, \beta')$ at $\beta' = \beta$ uses that the maximizer in $\widetilde{E}$ is $i$, giving $\partial V / \partial \beta = q - \PV^i[\1_S]$. Value maximization implies this derivative vanishes
\end{frame}

%--------------------------------------------------------------------
\begin{frame}\frametitle{Modigliani--Miller in the BGR framework}
\begin{proposition}[Modigliani--Miller in BGR]
Under the previous setup, if in addition every investor is either an equity holder or a bond holder, $I = I^e \cup I^b$, then for every $\beta'$ with $f_s(k) > \beta' \geq 0$ for all $s \in S$:
\begin{equation*}
V(k, \beta') = V(k, 0) = \widetilde{E}(k, 0) - k
\end{equation*}
\end{proposition}
\bigskip

The value map $V(k, \cdot)$ is independent of the financial policy on the no-default interior
\end{frame}

%--------------------------------------------------------------------
\begin{frame}\frametitle{When Modigliani--Miller fails}
\textbf{Case 1: Binding no-default constraint}
\begin{itemize}
\item If $\inf_s f_s(k) = \beta$, the firm sits on the boundary of $\mathcal{A}$
\item The marginal envelope condition no longer applies
\end{itemize}
\bigskip

\textbf{Case 2: Investors who are neither equity nor bond holders}
\begin{itemize}
\item There exist $i$ with $\PV^i[\1_S] < q$ \emph{and} $\PV^i[f_1(k) - \beta \1_S] < E$
\item Such investors are strictly infra-marginal in both markets
\item A change in $\beta$ may shift the identity of the marginal equity holder, changing $\widetilde{E}$ by more than $q \cdot d\beta$
\end{itemize}
\bigskip

Whether MM holds is now an \emph{equilibrium} question, answered by whether $I = I^e \cup I^b$
\end{frame}

%====================================================================
\section{Intermediated Short Sales}
%====================================================================

\begin{frame}\frametitle{Can intermediation restore efficiency?}
The previous sections took short-sale constraints as exogenous
\bigskip

\textbf{Question:} Can financial intermediaries who issue contracts replicating short positions restore the full unanimity and efficiency results?
\bigskip

A \emph{financial intermediary} is a risk-neutral, zero-profit institution that:
\begin{itemize}
\item issues long positions on the firm's equity to some investors
\item issues short positions on the firm's equity to other investors
\item the two positions net out except for a default friction
\end{itemize}
\end{frame}

%--------------------------------------------------------------------
\begin{frame}\frametitle{The default technology}
A simple default technology: in every state $s$, a fraction $\delta \in (0, 1)$ of the short position fails to be repaid
\bigskip

To insure his solvency, the intermediary holds $\gamma$ units of the firm's equity as collateral, satisfying
\begin{equation*}
m \leq m(1 - \delta) + \gamma, \quad \text{i.e.,} \quad \gamma \geq m \delta
\end{equation*}
where $m$ is the volume of intermediated long (= short) positions
\end{frame}

%--------------------------------------------------------------------
\begin{frame}\frametitle{The intermediary's problem}
Long and short positions trade at different prices $q^+$ and $q^-$
\bigskip

The intermediary solves
\begin{equation*}
\max \bigl\{ (q^+ - q^-) m - q \gamma \colon m \geq 0, \ \gamma \geq m \delta \bigr\}
\end{equation*}
\bigskip

\textbf{Constant returns to scale:} a positive $m$ at equilibrium requires zero profits, hence
\begin{equation*}
\gamma = m \delta \et q^+ - q^- = q \delta
\end{equation*}
\bigskip

The spread between long and short prices equals the per-unit collateral cost
\end{frame}

%--------------------------------------------------------------------
\begin{frame}\frametitle{Investor's budget set with intermediation}
Each investor holds long position $\phi^i_+ \geq 0$ and short position $\phi^i_- \geq 0$ on the intermediated derivative
\bigskip

\begin{block}{At $t = 0$}
\begin{small}
\begin{equation*}
c^i_0 + q \theta^i + q^+ \phi^i_+ + E \alpha^i \leq \omega^i_0 + (E + q \beta - k) \xi^i + q^- \phi^i_-
\end{equation*}
\end{small}
\end{block}

\begin{block}{At $t = 1$ in state $s$}
\begin{small}
\begin{equation*}
c^i_s + (1 - \delta)[f_s(k) - \beta] \phi^i_- \leq \omega^i_s + \theta^i + [f_s(k) - \beta](\alpha^i + \phi^i_+)
\end{equation*}
\end{small}
\end{block}
\bigskip

A long position substitutes direct equity ownership; a short position replicates a negative equity holding up to the default friction
\end{frame}

%--------------------------------------------------------------------
\begin{frame}\frametitle{Adjusted price perceptions}
When intermediation is active ($m > 0$), the Makowski map must reflect the value of equity as collateral:
\begin{equation*}
\widetilde{E}(k', \beta') = \max \biggl\{ \widetilde{E}^c(k', \beta'), \ \frac{\widetilde{q}^+(k', \beta') - \widetilde{q}^-(k', \beta')}{\delta} \biggr\}
\end{equation*}
\bigskip

where
\begin{itemize}
\item $\widetilde{E}^c(k', \beta') = \max_i \PV^i[f_1(k') - \beta' \1_S]$ is the consumer-only valuation
\item the second argument is the value of equity as collateral input to the intermediation technology
\end{itemize}
\end{frame}

%--------------------------------------------------------------------
\begin{frame}\frametitle{Two equilibrium regimes}
At equilibrium with $m > 0$, two regimes may occur:
\bigskip

\begin{block}{Regime $q > q^+$}
\begin{itemize}
\item Equity trades at a premium over the long position
\item All outstanding shares purchased by the intermediary, $\gamma = 1$
\item The premium $q - q^+$ compensates the intermediation collateral cost
\end{itemize}
\end{block}

\begin{block}{Regime $q = q^+$}
\begin{itemize}
\item Investors are indifferent between equity and the long position
\item Some shares are held by consumers, some by the intermediary
\end{itemize}
\end{block}
\end{frame}

%--------------------------------------------------------------------
\begin{frame}\frametitle{Unanimity and efficiency with intermediation}
\begin{proposition}[Unanimity and efficiency with intermediation]
At a competitive equilibrium of the economy with intermediated short sales:
\begin{itemize}
\item equity holders unanimously support the production and financial decisions of the firm
\item the equilibrium allocation is constrained Pareto optimal relative to the admissibility notion adapted to the enriched asset structure
\end{itemize}
\end{proposition}
\bigskip

The proof parallels the previous theorems with two modifications: the price perception is replaced by the adjusted version with collateral value, and admissibility is enlarged to include the derivative positions
\end{frame}

%====================================================================
\section{Bibliography}
%====================================================================

\begin{frame}\frametitle{The Bisin--Gottardi--Ruta framework}
\begin{thebibliography}{8}
\begin{small}
\beamertemplatearticlebibitems
    \bibitem{BisinGottardiRuta2014}
    Bisin, A., Gottardi, P. and Ruta, G.
    \newblock Equilibrium corporate finance and intermediation.
    \newblock {\em NBER Working Paper} 20345, 2014.
    \bibitem{BisinClementiGottardi2025}
    Bisin, A., Clementi, G. L. and Gottardi, P.
    \newblock Capital structure and hedging demand with incomplete markets.
    \newblock {\em Economic Journal}, 2025.
\end{small}
\end{thebibliography}
\bigskip

The first paper develops the framework of this part; the second is a recently published refinement
\end{frame}

%--------------------------------------------------------------------
\begin{frame}\frametitle{Makowski}
\begin{thebibliography}{8}
\begin{small}
\beamertemplatearticlebibitems
    \bibitem{Makowski1980}
    Makowski, L.
    \newblock A characterization of perfectly competitive economies with production.
    \newblock {\em Journal of Economic Theory} 22, 1980.
    \bibitem{Makowski1983a}
    Makowski, L.
    \newblock Competition and unanimity revisited.
    \newblock {\em American Economic Review} 73, 1983.
    \bibitem{Makowski1983b}
    Makowski, L.
    \newblock Competitive stock markets.
    \newblock {\em Review of Economic Studies} 50, 1983.
\end{small}
\end{thebibliography}
\end{frame}

%--------------------------------------------------------------------
\begin{frame}\frametitle{Related literature}
\begin{thebibliography}{8}
\begin{small}
\beamertemplatearticlebibitems
    \bibitem{AllenGale1991}
    Allen, F. and Gale, D.
    \newblock Arbitrage, short sales, and financial innovation.
    \newblock {\em Econometrica} 59, 1991.
    \bibitem{CarcelesCoenPirani2009}
    Carceles-Poveda, E. and Coen-Pirani, D.
    \newblock Shareholders' unanimity with incomplete markets.
    \newblock {\em International Economic Review} 50, 2009.
    \bibitem{AcharyaBisin2009}
    Acharya, V. V. and Bisin, A.
    \newblock Managerial hedging, equity ownership, and firm value.
    \newblock {\em RAND Journal of Economics} 40, 2009.
    \bibitem{BisinGottardiRampini2008}
    Bisin, A., Gottardi, P. and Rampini, A.
    \newblock Managerial hedging and portfolio monitoring.
    \newblock {\em Journal of the European Economic Association} 6, 2008.
\end{small}
\end{thebibliography}
\end{frame}

\end{document}